% ==============================================================================
%  4.6  Soluciones cuando el modelo no es válido
% ==============================================================================
\section{Soluciones cuando el modelo no es válido}
\label{sec:t4-soluciones}

Cuando el modelo no es válido hay, básicamente, dos opciones:
\begin{itemize}
  \item utilizar una metodología más apropiada, que con frecuencia será más complicada;
  \item encontrar una transformación de los datos para la que el modelo de regresión lineal sea
    adecuado.
\end{itemize}
Un modelo ajustado con alguna de las variables transformadas es un modelo nuevo, que requiere
el cumplimiento de las hipótesis: también es necesario validar los modelos transformados.

Según la hipótesis que falle, las soluciones habituales son las siguientes.
\begin{itemize}
  \item \textbf{Relación no lineal}:
    \begin{itemize}
      \item transformar $Y$, $X$ o ambas, de forma que se linealice la relación entre ellas;
      \item introducir términos polinómicos de orden mayor o nuevas variables independientes.
        Por ejemplo, podría ajustarse una función de regresión cuadrática,
        \[
          Y_i = \beta_0+\beta_1X_i+\beta_2X_i^2+\varepsilon_i ,
        \]
        o exponencial,
        \[
          Y_i = \beta_0\beta_1^{X_i}+\varepsilon_i .
        \]
    \end{itemize}
  \item \textbf{Varianza no constante}:
    \begin{itemize}
      \item transformar $Y$;
      \item utilizar otros métodos de ajuste, como el de mínimos cuadrados ponderados.
    \end{itemize}
  \item \textbf{Errores no normales}: transformar $Y$. Con frecuencia la falta de normalidad y la
    heterogeneidad de varianzas se presentan juntas; lo habitual es buscar primero una
    transformación que estabilice la varianza.
\end{itemize}

Muchas veces, transformaciones sencillas de $X$, de $Y$ o de ambas permiten que los modelos
lineales resulten adecuados con las variables transformadas.

\subsection{Transformaciones para linealizar relaciones no lineales}

Supongamos que es razonable considerar que los errores son normales y con varianza constante.
\begin{itemize}
  \item Las transformaciones de $Y$ podrían modificar la forma de su distribución y, por tanto,
    la de los errores, incluidas la normalidad y la homogeneidad de varianzas.
  \item Las transformaciones de $X$ pueden ayudar a linealizar la relación sin afectar a la
    distribución de $Y$.
\end{itemize}
En este caso conviene, por tanto, transformar $X$. Las transformaciones adecuadas dependen de la
forma de la relación observada en el diagrama de dispersión:
\begin{center}
  \renewcommand{\arraystretch}{1.3}
  \begin{tabular}{ll}
    \toprule
    Patrón de la relación & Transformaciones de $X$ \\
    \midrule
    Creciente, cada vez más lentamente (cóncava) & $X'=\log_{10}X$, \quad $X'=\sqrt{X}$ \\
    Creciente, cada vez más rápidamente (convexa) & $X'=X^2$, \quad $X'=\exp(X)$ \\
    Decreciente, cada vez más lentamente (convexa) & $X'=1/X$, \quad $X'=\exp(-X)$ \\
    \bottomrule
  \end{tabular}
\end{center}

\begin{nota}
  Puede ser necesario incluir alguna constante en la transformación para que tenga sentido en
  todos los valores observados. Por ejemplo, $X'=\sqrt{X+k}$, con $k$ tal que $X+k\ge0$.
\end{nota}

\subsection{Transformaciones para la heterogeneidad de varianzas y la falta de normalidad}

Con frecuencia la falta de normalidad y la heterogeneidad de varianzas aparecen juntas. En ese
caso se transforma $Y$ para modificar la forma de su distribución, y a veces también se consigue
linealizar la relación. Una transformación muy utilizada es el \textbf{logaritmo}.

\begin{ejemplo}[Transformación logarítmica][ej:t4-log]
  \begin{itemize}
    \item El histograma de una variable de gastos es muy asimétrico a la derecha: la mayoría de
      los valores se concentran cerca de 0 y unos pocos son muy grandes. El histograma del
      logaritmo de los gastos es, en cambio, aproximadamente simétrico y con forma de campana.
    \item Al representar el salario de los empleados de una empresa frente a su antigüedad (en
      años), la relación es creciente y aproximadamente lineal, pero la dispersión de los
      salarios aumenta con la antigüedad. En el gráfico de residuos frente a valores ajustados se
      observa que la dispersión de los residuos crece con $\est{y}_i$ (forma de embudo). Si se
      ajusta el modelo con $\log(\text{salario})$ como respuesta, la relación con la antigüedad
      sigue siendo lineal y los residuos presentan una dispersión aproximadamente constante.
  \end{itemize}
\end{ejemplo}

Además del logaritmo, hay otras transformaciones útiles de $Y$, como la raíz cuadrada o la
inversa, que forman parte de la familia de transformaciones potencia. Para elegir la
transformación que mejor soluciona el problema se utiliza la transformación de Box--Cox, que
selecciona dentro de esa familia la más adecuada a los datos.

\begin{definicion}[Transformación de Box--Cox][def:t4-boxcox]
  La \textbf{transformación de Box--Cox} es la familia de transformaciones no lineales, para
  $y>0$,
  \[
    g_\lambda(y) =
    \begin{cases}
      \dfrac{y^\lambda-1}{\lambda} & \text{si } \lambda\neq0, \\[8pt]
      \log(y) & \text{si } \lambda=0 .
    \end{cases}
  \]
  Cuando $Y$ puede tomar valores negativos, se utiliza
  \[
    g_\lambda(y) =
    \begin{cases}
      \dfrac{(y+k)^\lambda-1}{\lambda} & \text{si } \lambda\neq0, \\[8pt]
      \log(y+k) & \text{si } \lambda=0 ,
    \end{cases}
  \]
  con $k$ tal que $y+k>0$ para todos los valores observados.
\end{definicion}

\begin{itemize}
  \item La transformación depende del parámetro $\lambda$, que puede elegirse maximizando la
    verosimilitud, suponiendo que para algún $\lambda$ dado las variables transformadas siguen una
    distribución $\Normal(\mu,\sigma^2)$. En la práctica se representa la log-verosimilitud en
    función de $\lambda$ y se toma el valor que la maximiza, junto con un intervalo de confianza
    del 95\,\% para $\lambda$. Por ejemplo, si el máximo se alcanza en $\lambda\approx0$ y el
    intervalo es aproximadamente $(-0.5,\,0.5)$, se elige la transformación logarítmica.
  \item $\lambda=1$ corresponde a la transformación $y'=y-1$, que solo supone un cambio en el
    \emph{intercept} del modelo y no en la forma de la distribución; es decir, equivale a no
    transformar.
\end{itemize}
